%!TEX root = ../main.tex

\chapter*{\label{ch:conclusions}Conclusions}
\markboth{Conclusions}{}
\addcontentsline{toc}{chapter}{Conclusions} 
% NOTE: Introduction and Conclusions are special chapters: they are not
%       marked by a section number, so we need to add it manually to the
%       ToC, as done above.

Building upon the theoretical foundations of neuroscientific models (Chapter 1) and Scientific Machine Learning (Chapter 2), Chapter 3 demonstrated the expressive fitting power of Operator Learning on both synthetic and empirical data. Regarding the synthetic approach, an interesting avenue for future research would be to determine the maximum network size, $N$, up to which the performance of our setup remains stable. Conversely, concerning real-world applications, future efforts should focus on refining neural operator architectures to achieve the quantitative precision required for clinical translation. Specifically, minimizing this error profile will enable the development of model-free anomaly detection systems. Trained exclusively on baseline data from healthy subjects, these refined operators could identify subtle, statistically significant deviations in neural output, thereby serving as non-invasive diagnostic tools for neurological disorders. Furthermore, the generalizability of these models warrants testing across adjacent domains, including computational social science and bioinformatics.

Finally, Chapter 4 demonstrated that expressive power and interpretability can successfully coexist. In this context, an important benchmark omitted from this study is the Graph Neural Network (GNN), which naturally processes network structures via message-passing operations. It would be valuable to evaluate whether the LFNO outperforms the GNN when both are deployed within our framework. Furthermore, due to time constraints, the performance of the LFNO has yet to be benchmarked against empirical datasets. A key milestone for future research will be validating the LFNO framework on large-scale real-world data, such as modeling epidemic propagation across the complex network of Italian regional connections.

Crucially, the most promising and mathematically profound extension of this work lies in the reverse-engineering of underlying physical laws directly from observational data (Section \ref{sec:Inverse Problem}). By exploiting the trained LFNO, we reconstructed the reaction potentials of interacting species, mapping the data onto an empirical hypersurfaces within the system's phase space. Through the subsequent application of symbolic regression, we anticipate distilling these learned hypersurfaces into explicit, closed-form analytical expressions derived purely from data. In an epidemiological context, this translates to discovering the exact functional forms of transmission and recovery rates directly from noisy surveillance data, without relying on prior mechanistic assumptions. Beyond this inverse problem, future projects will explore inferring hidden network topologies when reaction terms are known, quantifying the model's robustness against topological uncertainties, and extending the LFNO mathematical framework from undirected graphs to complex directed networks.